\documentclass[11pt,a4paper]{article}

% --- Paquetes Básicos y de Formato ---
\usepackage[utf8]{inputenc}
\usepackage[spanish,es-tabla]{babel} % Configura el idioma español y cambia "Cuadro" por "Tabla" si se desea, o mantiene el estándar.
\usepackage[margin=2.5cm,top=3cm,bottom=2.5cm]{geometry} % Márgenes institucionales estándar
\usepackage{graphicx}
\usepackage{booktabs} % Para tablas con acabado profesional
\usepackage{tabularx} % Para que las tablas se ajusten al ancho de la página
\usepackage{titlesec} % Para personalizar los títulos
\usepackage{parskip}  % Añade espacio entre párrafos automáticamente sin sangría

% --- Configuración de Títulos ---
\titleformat{\section}{\normalfont\Large\bfseries}{\thesection.}{0.5em}{}
\titleformat{\subsection}{\normalfont\large\bfseries}{\thesubsection.}{0.5em}{}

\begin{document}

% --- Encabezado Oficial ---
\begin{center}
    \small\MakeUppercase{“Año de la recuperación y consolidación de la economía peruana”} [cite: 1] \\
    \vspace{0.8cm}
    \Large\bfseries INFORME 001-2025-APAMNRN/COORDINADOR-PNT/LHQ [cite: 2]
\end{center}

\vspace{0.6cm}

% --- Bloque de Datos del Informe ---
\begin{tabular}{@{}llp{11cm}@{}}
    \bfseries A & : & SR. MARCELO DIAZ HUAHUASONCCO \\ [cite: 3]
                &   & PRESIDENTE DE LA APA MICROCUENCA NORTE RIO NUÑOA \\ [cite: 4] \vspace{0.15cm} \\
    \bfseries DE & : & MVZ. LUIS HUAMAN QUISPE \\
                 &   & COORDINADOR DEL PNT \\ [cite: 6] \vspace{0.15cm} \\
    \bfseries ASUNTO & : & INFORME DE ACTIVIDADES MES DE NOVIEMBRE \\ [cite: 7] \vspace{0.15cm} \\
    \bfseries REFERENCIA & : & CONTRATO Nº005-2025 \\ [cite: 8] \vspace{0.15cm} \\
    \bfseries FECHA & : & Nuñoa, 15 de Diciembre del 2025. [cite: 9]
\end{tabular}

\vspace{0.3cm}
\noindent\rule{\textwidth}{1pt} % Línea divisoria horizontal
\vspace{0.4cm}

% --- Introducción ---
Previo cordial saludo, mediante el presente documento es grato dirigirme a Ud. [cite: 10] para hacer alcance del informe mensual correspondiente al mes de noviembre de las actividades realizadas como coordinador del PNT “MEJORAMIENTO DEL MANEJO PRODUCTIVO DE LA ALPACA” en la ASOCIACION DE PRODUCTORES AGROPECUARIOS MICROCUENCA NORTE RIO NUÑOA, según documento de Referencia, adjuntando también otros documentos al presente. [cite: 11]

\vspace{0.2cm}

% --- Sección I ---
\section{INFORMACION GENERAL DEL PNT.} [cite: 12]

CUADRO Nº1. Información general del Plan de Negocio. [cite: 13]

\begin{table}[h!]
    \centering
    \small
    \begin{tabularx}{\textwidth}{|l|X|}
        \hline
        \bfseries NOMBRE DE LA OA & ASOCIACIÓN DE PRODUCTORES AGROPECUARIOS MICROCUENCA NORTE RIO NUÑOA \\ \hline [cite: 14]
        \bfseries PLAN DE NEGOCIO & MEJORAMIENTO DEL MANEJO PRODUCTIVO DE LA ALPACA \\ \hline [cite: 14]
        \bfseries Nº DE SOCIOS & 27 \\ \hline [cite: 14]
        \bfseries COMUNIDAD & MICROCUENCA NORTE RIO NUÑOA \\ \hline [cite: 14]
        \bfseries DISTRITO & NUÑOA \\ \hline [cite: 14]
    \end{tabularx}
\end{table}

% --- Sección II ---
\section{ACTIVIDADES REALIZADAS DURANTE EL MES} [cite: 15]

\subsection{EVALUACIÓN Y SELECCIÓN DE REPRODUCTORES (PRE-EMPADRE)} [cite: 16]
Ante la proximidad de la época de empadre (Enero-Marzo), se realizó la asistencia técnica para la selección de los reproductores machos y hembras del plantel. [cite: 17]

\textbf{Evaluación Andrológica (Machos):} Se revisó la conformación física, testículos (tamaño y consistencia) y la liberación del pene (ausencia de adherencias) para asegurar la capacidad reproductiva. [cite: 18]

\textbf{Selección de Hembras:} Se identificaron las hembras aptas para el servicio, segregando aquellas con problemas reproductivos o mala condición corporal, optimizando así los recursos forrajeros disponibles. [cite: 19]

\subsection{CAPACITACIÓN: SANIDAD ANIMAL EN ÉPOCA DE LLUVIAS} [cite: 20]
Se llevó a cabo la capacitación técnica dirigida a prevenir las pérdidas económicas causadas por el clima adverso. [cite: 21]

\textbf{Temas Tratados:} Identificación de síntomas de neumonía en crías y adultos, y otras potenciales enfermedades estacionarias. [cite: 22]

\textbf{Resultado:} Los productores conocen las medidas preventivas de manejo sanitario para mitigar el impacto de las lluvias en la salud del rebaño. [cite: 23]

\subsection{GEORREFERENCIACIÓN PRELIMINAR PARA LA INSTALACIÓN DE LOS CERCOS ALPAQUEROS} [cite: 24]
Se procedió en el trabajo de ubicación y trazado de los módulos de recuperación de praderas nativas (cercos). [cite: 25]

\textbf{Actividad:} Identificación de los vértices donde se instalarán los postes metálicos para los cercos de mallas ganadera. [cite: 26]

\textbf{Marcación:} Se realizó el estacado y marcado con yeso de los puntos esquineros para definir el perímetro de los módulos de una hectárea, optimizando la distribución de los cercos. [cite: 27]

% --- Sección III ---
\section{METODOLOGÍA Y RECURSOS UTILIZADOS} [cite: 28]

\subsection{METODOLOGÍA} [cite: 29]
\begin{itemize}
    \item \textbf{Capacitación:} Método demostrativo participativo en el centro de manejo. [cite: 30]
    \item \textbf{Asistencia Técnica:} Visitas personalizadas a predios para la verificación de suelos y trazado de áreas. [cite: 31]
\end{itemize}

\subsection{RECURSOS UTILIZADOS} [cite: 32]
\begin{itemize}
    \item Motocicleta para desplazamiento. [cite: 33]
    \item Equipo de esquila mecánica. [cite: 34]
    \item Materiales de campo: wincha. [cite: 35]
    \item Computadora. [cite: 36]
    \item Registro de asistencia. [cite: 37]
    \item Otros. [cite: 38]
\end{itemize}

% --- Sección IV ---
\section{CONCLUSIONES} [cite: 39]
\begin{itemize}
    \item Los productores han sido capacitados en medidas preventivas sanitarias, estando mejor preparados para enfrentar la temporada crítica de lluvias y reducir la mortalidad por enfermedades infecciosas. [cite: 40]
    \item Se ha realizado la selección técnica de los reproductores, garantizando que solo los animales aptos ingresen a la campaña de empadre de enero. [cite: 41]
    \item Existe un avance físico en la preparación del terreno para la instalación de cercos, optimizando el tiempo y la mano de obra disponible. [cite: 42]
\end{itemize}

% --- Sección V ---
\section{RECOMENDACIONES} [cite: 43]
\begin{itemize}
    \item \textbf{A los Socios:} Poner en práctica las recomendaciones de manejo sanitario para evitar la aparición de enfermedades en crías y adultos por la temporada de lluvias próxima. [cite: 44]
    \item \textbf{A la Directiva:} Agilizar la gestión de compra de los insumos restantes para iniciar la dosificación y el cerrado de cercos a la brevedad. [cite: 45]
\end{itemize}

% --- Sección VI ---
\section{ANEXOS} [cite: 46]
\begin{itemize}
    \item \textbf{Anexo 01:} Registro de asistencia al taller de capacitación. [cite: 47]
    \item \textbf{Anexo 02:} Panel fotográfico (Capacitación, Selección de animales, Cercos). [cite: 48]
    \item \textbf{Anexo 03:} Fichas de asistencia técnica. [cite: 49]
    \item \textbf{Anexo 04:} Copia de Recibo por Honorarios. [cite: 50]
\end{itemize}

\newpage % Salto de página para los informes de talleres adjuntos

% =========================================================================
% INFORME ADJUNTO: TALLER DE CAPACITACIÓN 1
% =========================================================================
\begin{center}
    \Large\bfseries TALLER DE CAPACITACIÓN [cite: 51] \\
    \large TEMA: PREVENCIÓN DE ENFERMEDADES EN ALPACAS DURANTE LA TEMPORADA DE LLUVIAS [cite: 52]
\end{center}

\vspace{0.4cm}

\textbf{OBJETIVO} [cite: 53]
\begin{itemize}
    \item Que los participantes identifiquen los principales riesgos sanitarios asociados al clima adverso (lluvias, heladas y granizadas), centrándose específicamente en la prevención de neumonías y enterotoxemias. [cite: 54]
    \item Capacitar a los productores en el manejo preventivo de la infraestructura (cobertizos y dormideros) y el uso correcto del botiquín veterinario de emergencia para reducir la mortalidad en el hato. [cite: 55]
\end{itemize}

\textbf{DESARROLLO DEL EVENTO} [cite: 56]

\textbf{Fase Teórica (Introducción):} Se inició la sesión analizando cómo el exceso de humedad y los cambios bruscos de temperatura afectan las defensas de las alpacas, especialmente en crías y madres gestantes. [cite: 57] El ponente explicó los síntomas clínicos de la \textit{Neumonía} (tos, fiebre, dificultad respiratoria) y la \textit{Enterotoxemia} (muerte súbita, diarrea), enfatizando que la prevención es más económica que el tratamiento. [cite: 58]

\textbf{Demostración Práctica "In Situ":} Se procedió a la revisión de los dormideros y cobertizos. [cite: 59] Se demostró cómo identificar puntos de encharcamiento y la importancia de mantener la "cama" (suelo) seca utilizando paja o guano seco para evitar el enfriamiento nocturno de los animales. [cite: 60] Asimismo, se realizó una revisión del \textit{Botiquín Veterinario}, enseñando a identificar los antibióticos, vitaminas y reconstituyentes necesarios para esta época. [cite: 61]

\textbf{Participación de Productores:} Los productores participaron activamente identificando las mejoras necesarias en sus corrales para mejorar el drenaje. [cite: 62] Además, practicaron el cálculo de dosis de medicamentos preventivos (vitaminas) según el peso vivo del animal (al ojo) bajo la supervisión del técnico. [cite: 63]

\textbf{Cierre:} Se concluyó con una ronda de preguntas sobre el cuidado de las hembras en el último tercio de gestación y la importancia de la colostración temprana en las crías que nacerán pronto. [cite: 64] Se programó la siguiente visita para el monitoreo de parición. [cite: 65]

\vspace{0.3cm}
\textbf{LUGAR Y FECHA} [cite: 66]
\begin{tabular}{@{}lll@{}}
    \textbf{Provincia:}   & Melgar [cite: 67] \\
    \textbf{Distrito:}    & Nuñoa [cite: 68] \\
    \textbf{Lugar:}       & Local de la Asociación [cite: 69] \\
    \textbf{Fecha y hora:} & 10 de Diciembre, 10:00 am [cite: 70]
\end{tabular}

\vspace{0.3cm}
\textbf{RESULTADOS} [cite: 71]
\begin{itemize}
    \item Los productores han adquirido conocimientos para identificar tempranamente los signos de enfermedades respiratorias y entéricas, permitiendo una reacción rápida para salvar al animal. [cite: 72]
    \item Los participantes reconocen la importancia de mantener los dormideros secos y limpios como la principal medida para prevenir la mortalidad durante la temporada de lluvias. [cite: 73]
\end{itemize}

\textbf{ANEXOS} [cite: 74]
\begin{itemize}
    \item Registro de participantes. [cite: 75]
    \item Registro fotográfico. [cite: 76]
\end{itemize}

\newpage

% =========================================================================
% INFORME ADJUNTO: TALLER DE CAPACITACIÓN 2
% =========================================================================
\begin{center}
    \Large\bfseries TALLER DE CAPACITACIÓN [cite: 77] \\
    \large TEMA: TÉCNICAS DE ESQUILA MECÁNICA EN ALPACAS Y MANEJO DE VELLÓN [cite: 78]
\end{center}

\vspace{0.4cm}

\textbf{OBJETIVO} [cite: 79]
\begin{itemize}
    \item Que los participantes conozcan y se familiaricen con el funcionamiento, armado y mantenimiento de los equipos de esquila mecánica, diferenciando las ventajas frente a la esquila manual (tijera). [cite: 80]
    \item Realizar la demostración práctica "in situ" del proceso de esquila, asegurando que los productores adopten técnicas que garanticen el bienestar animal y la obtención de un vellón de calidad comercial. [cite: 81]
\end{itemize}

\textbf{DESARROLLO DEL EVENTO} [cite: 82]

\textbf{Fase Teórica (Introducción):} Se inició la sesión compartiendo saberes previos sobre la época de esquila y sus métodos tradicionales. [cite: 83] El ponente explicó las partes del equipo de esquila mecánica (motor, manga, pieza de mano, peines y cortantes) y las normas de seguridad industrial para evitar accidentes durante la operación. [cite: 84]

\textbf{Demostración Práctica "In Situ":} Se procedió a la práctica en campo con los animales. [cite: 85] Se demostró la correcta sujeción (tumbado y manea) para evitar el estrés del animal (bienestar animal). [cite: 86] El ponente realizó la esquila demostrativa explicando las posiciones correctas para evitar cortes en la piel y "dobles cortes" en la fibra. [cite: 87]

\textbf{Participación de Productores:} Los productores participaron activamente bajo supervisión, manipulando la esquiladora y practicando el desbordado (envellonado) básico de la fibra para separar impurezas. [cite: 88]

\textbf{Cierre:} Se concluyó con una ronda de preguntas sobre el afilado de peines y el mantenimiento post-esquila de la maquinaria. [cite: 89] Se programaron las fechas para la siguiente asistencia técnica. [cite: 90]

\vspace{0.3cm}
\textbf{LUGAR Y FECHA} [cite: 91]
\begin{tabular}{@{}lll@{}}
    \textbf{Provincia:}    & Melgar [cite: 92] \\
    \textbf{Distrito:}     & Nuñoa [cite: 93] \\
    \textbf{Lugar:}        & Fundo de uno de los socios participantes [cite: 94] \\
    \textbf{Fecha y hora:} & \line(1,0){40} % Línea vacía ya que no se especificó hora en el original [cite: 95]
\end{tabular}

\vspace{0.3cm}
\textbf{RESULTADOS} [cite: 96]
\begin{itemize}
    \item Los productores han adquirido conocimientos sobre el uso eficiente de la esquila mecánica, reduciendo el tiempo de manejo por animal y minimizando el estrés en el hato alpaquero. [cite: 97]
    \item Los participantes reconocen la importancia de un buen corte para no afectar el precio de la fibra y muestran interés en aplicar la categorización básica del vellón durante la cosecha. [cite: 98]
\end{itemize}

\textbf{ANEXOS} [cite: 99]
\begin{itemize}
    \item Registro de participantes. [cite: 100]
    \item Registro fotográfico. [cite: 101]
\end{itemize}

\end{document}
